
\section{What a representation-risk certificate establishes}
\label{sec:object}

Let $X$ denote the information available at a forecast origin and $Y$ its
scalar observed target. Condition on all development data and choices, so
the representation $\Gamma:\mathcal X\to\mathcal U$, its measurable
finite-valued pseudometric $d$, and a sensitivity budget $\kappa\geq0$ are
fixed. Assume $Y$ lies almost surely in a known interval
$I=[y_-,y_+]$ of diameter $D>0$, and put $B=D^2$.
Define
\begin{align}
 \mathcal G_\kappa(d)
 &=\{g:\mathcal U\to I\text{ measurable}:
       |g(u)-g(v)|\leq\kappa d(u,v)\ \text{for all }u,v\},
       \label{eq:class}\\
 R(g)&=\mathbb E[(Y-g(\Gamma(X)))^2],
 &R^*_\kappa&=\inf_{g\in\mathcal G_\kappa(d)}R(g).
       \label{eq:risk}
\end{align}
All displayed infima and statistics are assumed measurable. The expectation
uses the declared evaluation law, conditional on development, and not the
law conditional on an individual query. Scalar clipping to $I$ preserves
the Lipschitz constraint and decreases squared loss. Hence this bounded
class has the same optimal risk as the corresponding unrestricted scalar
Lipschitz class. Constants belong to it, giving
$0\leq R^*_\kappa\leq\operatorname{Var}(Y)\leq B/4$.
When $D=0$, the target and interval-valued predictions are the same constant;
all risks and certificates are zero, without any division by $B$.

The representation, metric, budget, target, loss, and evaluation law jointly
specify the certification problem. Changing feature scaling while retaining
the same numerical budget generally changes the class. A ridge penalty
does not establish a hard Lipschitz budget. A fitted upper endpoint can be
paired with a class lower endpoint only after establishing that the fitted
decoder $\widehat g$ belongs to that class.

The standard conditional-expectation decomposition explains the scope of
this risk. With $M_X=\mathbb E[Y\mid X]$ and
$M_U=\mathbb E[Y\mid\Gamma(X)]$, orthogonality in $L^2$ gives
\begin{equation}
 R^*_\kappa
 =\underbrace{\mathbb E(Y-M_X)^2}_{\text{future uncertainty}}
  +\underbrace{\mathbb E(M_X-M_U)^2}_{\text{conditional-mean information loss}}
  +\underbrace{\inf_{g\in\mathcal G_\kappa(d)}
       \mathbb E(M_U-g(\Gamma(X)))^2}_{\text{sensitivity restriction}}.
 \label{eq:decomposition}
\end{equation}
The three residual components are pairwise orthogonal because
$\Gamma(X)$ is measurable with respect to $X$; taking the infimum yields
the identity. A positive lower risk certificate does not separate these
terms. It is therefore neither a certificate of unrestricted information
loss nor a sufficient rule for discovering a successful refinement.

\subsection{The pairwise starting point}

Draw a variable $A$, and then draw $(U,Y)$ and $(U',Y')$ independently from
the same conditional law given $A$, where each copy has the declared
marginal law and $U=\Gamma(X)$. With $z_+=\max\{z,0\}$, define
\begin{equation}
 W_{\kappa,A}
 =\frac12\mathbb E\bigl(|Y-Y'|-\kappa d(U,U')\bigr)_+^2.
 \label{eq:pair}
\end{equation}
Taking $A$ constant gives unconditional independent pairing. A common
forecast origin or a frozen stratum gives a different admissible
construction when its conditional independence and marginal-law conditions
hold. Arbitrary nearest-neighbor matching need not meet these conditions.

\begin{proposition}[Pairwise obstruction]
\label{prop:pair}
Under this pairing construction, $W_{\kappa,A}\leq R^*_\kappa$.
\end{proposition}

\begin{proof}
For any fixed admissible $g$, put $r=Y-g(U)$ and $r'=Y'-g(U')$.
The triangle inequality implies
$(|Y-Y'|-\kappa d(U,U'))_+\leq|r-r'|$.
Conditional independence gives
$\frac12\mathbb E(r-r')^2
=\mathbb E r^2-\mathbb E\{\mathbb E[r\mid A]^2\}\leq R(g)$.
Square the pointwise inequality, take expectations, and then take the
infimum over $g$.
\end{proof}

Paired Lipschitz obstruction arguments have prior precedents
\citep{akash2021inverse}. The role of the pair score here is its exact
connection to a centered hierarchy, including explicit cases in which
pairwise information is arbitrarily loose.

\section{A centered hierarchy and its finite-block limits}
\label{sec:hierarchy}

For $b\geq2$ unconditional iid observations $(U_i,Y_i)$ from the fixed
evaluation law, define
\begin{equation}
 T_b=\frac1{b-1}\min_{v\in\mathbb R^b}
 \left\{\sum_{i=1}^b(Y_i-v_i)^2:
    |v_i-v_j|\leq\kappa d(U_i,U_j)\quad(i,j\leq b)\right\},
 \qquad \mathcal W_b=\mathbb E T_b.
 \label{eq:block}
\end{equation}
The feasible set is nonempty and closed and the objective is strictly
convex and coercive, so the minimizer exists and is unique. Clipping its
values to $I$ preserves feasibility and cannot increase the objective.
If $a_b$ is the usual empirical minimum MSE, then
$T_b=b a_b/(b-1)$. The factor $1/(b-1)$ exploits translation-invariant
difference constraints; it is not a generic normalization for empirical
risk minimization.

\begin{theorem}[Centered scalar block hierarchy]
\label{thm:hierarchy}
For the fixed class in \eqref{eq:class} and bounded scalar target above:
\begin{enumerate}
\item $T_2=\frac12(|Y_1-Y_2|-\kappa d(U_1,U_2))_+^2$, so
$\mathcal W_2=W_{\kappa,A}$ for constant $A$.
\item With
\begin{equation}
 C_b=\frac{B\lfloor b^2/4\rfloor}{b(b-1)},
 \qquad 0\leq T_b\leq C_b,
 \qquad \mathcal W_b\leq R^*_\kappa.
 \label{eq:range}
\end{equation}
\item If $T_m(S)$ denotes the score on a subset $S$ of $m$ observations
from a realized $b$-block, then for $2\leq m\leq b$,
\begin{equation}
 \binom bm^{-1}\sum_{|S|=m}T_m(S)\leq T_b.
 \label{eq:subset}
\end{equation}
Consequently $\mathcal W_2\leq\mathcal W_3\leq\cdots\leq R^*_\kappa$.
\item If $(\mathcal U,d)$ is a compact metric space, then
$\mathcal W_b\uparrow R^*_\kappa$ as $b\to\infty$.
\end{enumerate}
\end{theorem}

\begin{proof}
For two observations, if the target separation exceeds the allowed fitted
separation by $e>0$, the optimum has opposite residuals $e/2$ and $-e/2$.
Its squared residual sum is $e^2/2$. Otherwise both targets are feasible.
This proves the first assertion.

For any fixed admissible $g$, let $r_i=Y_i-g(U_i)$ and
$\bar r=b^{-1}\sum_i r_i$. The translated fitted values
$g(U_i)+\bar r$ satisfy every difference constraint. Thus
\[
 T_b\leq\frac1{b-1}\sum_i(r_i-\bar r)^2,
 \qquad \mathcal W_b\leq\operatorname{Var}(Y-g(U))\leq R(g).
\]
Taking the infimum gives the risk bound. Choosing the constant fitted value
$\bar Y$ bounds $T_b$ by the unbiased sample variance of $Y_1,\ldots,Y_b$.
That variance is convex in each target coordinate; its maximum on $I^b$
occurs at interval endpoints. With $k$ upper endpoints its value is
$B k(b-k)/\{b(b-1)\}$. Maximizing over integer $k$ proves the range.

Let $v$ minimize the full-block objective and set $r_i=Y_i-v_i$.
Translation invariance implies $\sum_i r_i=0$. Therefore
\[
 T_b=\frac1{b(b-1)}\sum_{i<j}(r_i-r_j)^2.
\]
Restricting $v$ to $S$ and translating by that subset's residual mean gives
\[
 T_m(S)\leq\frac1{m(m-1)}
              \sum_{i<j,\ i,j\in S}(r_i-r_j)^2.
\]
A uniformly chosen $m$-subset contains each pair with probability
$m(m-1)/\{b(b-1)\}$. Averaging yields \eqref{eq:subset}; taking
expectations gives monotonicity. This does not order arbitrary realized
blocks of different sizes.

For consistency, every finite feasible vector admits the scalar extension
$g_v(u)=\min_i\{v_i+\kappa d(u,U_i)\}$. The triangle inequality makes it
$\kappa$-Lipschitz and the finite constraints give $g_v(U_i)=v_i$.
After clipping to $I$, this identifies $a_b$ with the empirical infimum over
the same bounded Lipschitz class as the population problem. On a compact
metric space that class is compact in the uniform norm by boundedness and
equicontinuity. Choose a finite uniform $1/j$-net for each integer $j$.
The strong law holds simultaneously on their countable union. Since
\[
 |(y-f(u))^2-(y-g(u))^2|\leq2D\|f-g\|_\infty
 \quad(y,f(u),g(u)\in I),
\]
net approximation yields uniform convergence of empirical risks and hence
$a_b\to R^*_\kappa$ almost surely. Now $T_b=b a_b/(b-1)$ and
$0\leq T_b\leq B/2$. Dominated convergence proves the claimed limit.
\end{proof}

The strengthened information comes from requiring jointly consistent fitted
values across all pairs in a block. Monotonic expected sample-average
minima and batch lower estimation are established precedents
\citep{mak1999montecarlo,linderoth2006empirical}. The specific connection
here is the centered normalization, its exact pair member, its subset
inequality, and its Lipschitz-risk limit. It does not supply a new
concentration inequality or a vector-target consistency theorem.

\subsection{No fixed block size is uniformly informative}

\begin{proposition}[Rare-conflict and noisy-support laws]
\label{prop:limits}
For $(U,Y)=(3,0),(0,-1),(1,1)$ with probabilities
$1-2\varepsilon,\varepsilon,\varepsilon$, where
$0<\varepsilon<1/2$, and $\kappa=1$ with ordinary distance,
\begin{equation}
 R^*_\kappa=\varepsilon/2,
 \qquad
 \mathcal W_b=
 \frac{2b\varepsilon-1+(1-2\varepsilon)^b}{4(b-1)}.
 \label{eq:rare}
\end{equation}
For $U$ uniform on $M$ points with pairwise distances at least $D/\kappa$,
where $D,\kappa>0$, and an independent equiprobable target
$Y\in\{0,D\}$,
\begin{equation}
 R^*_\kappa=D^2/4,
 \qquad
 \mathcal W_b=\frac{D^2}{4(b-1)}
       \left[b-M\{1-(1-1/M)^b\}\right].
 \label{eq:occupancy}
\end{equation}
For each fixed $b$, the ratio $\mathcal W_b/R^*_\kappa$ tends to zero as
$\varepsilon\downarrow0$ in the first law or $M\to\infty$ in the second.
\end{proposition}

\begin{proof}
In the first law, the two rare fitted values must differ by at most one.
Their minimum unweighted sum of squared errors is $1/2$, attained at
$-1/2$ and $1/2$; the common atom can be fitted at zero. A scalar
Lipschitz extension realizes these values, proving the population optimum.
In a block, let $N_-,N_+$ count the rare atoms and put $N=N_-+N_+$.
Their constrained weighted residual sum is $N_-N_+/N$ when $N>0$ and
zero otherwise. Since $N\sim\operatorname{Binomial}(b,2\varepsilon)$ and
$N_-\mid N\sim\operatorname{Binomial}(N,1/2)$, the expected sum is
$\{\mathbb EN-\Pr(N\geq1)\}/4$. Dividing by $b-1$ proves
\eqref{eq:rare}.

In the second law, all interval-valued fitted values at distinct support
points satisfy the distance constraints. At a point observed $n_j>0$
times, fitting its target mean has expected residual sum
$(n_j-1)D^2/4$. If $O_b$ is the number of occupied points, the total
expected sum is $D^2\mathbb E(b-O_b)/4$, and
$\mathbb EO_b=M\{1-(1-1/M)^b\}$. The constant $D/2$ attains the
population risk $D^2/4$ because the target is independent of $U$.
The stated limits follow from the displayed formulas.
\end{proof}

The noisy-support examples are finite compact spaces; their diameter can
remain $D/\kappa$ as $M$ grows. Thus block size and target range, even
together with a bound on representation diameter, do not yield a uniform
vanishing approximation error over these spaces. Consistency is a statement
for a fixed evaluation law, not a universal finite-block guarantee.

\paragraph{Common forecast origins require a different limit.}
If the observations in a block are independent only conditional on a common
$A$, validity and monotonicity still hold. For finite-valued $A$ under the
corresponding compactness conditions, the limiting quantity is
\[
 \sum_a\Pr(A=a)\inf_{g\in\mathcal G_\kappa(d)}
                  \mathbb E[(Y-g(U))^2\mid A=a].
\]
It need not equal the risk of one shared optimal decoder. For example,
with constant $U$ and a target determined by a nonconstant common origin,
all within-origin block scores vanish although the shared constant decoder
has positive risk. Drawing an independent origin in each independent
episode instead produces unconditional iid observations from the declared
episode--origin mixture.

\section{Computable lower values and simultaneous calibration}
\label{sec:calibration}

Optimization validity and statistical calibration are separate steps. A
primal objective bounds an empirical minimum from above and therefore
cannot supply the lower value needed below. Sparse or unfinished dual
calculations can supply valid lower values if their arithmetic is verified.

\subsection{A deterministic finite-cohort certificate}

For any fixed cohort $(u_i,y_i)_{i=1}^n$ with $y_i\in I$, define
$a_n=\inf_{g\in\mathcal G_\kappa(d)}n^{-1}\sum_i(y_i-g(u_i))^2$.
The scalar extension used above identifies $n a_n$ with the finite
quadratic minimum under all pairwise Lipschitz constraints. This equality
uses no independence or population sampling assumption.

\begin{proposition}[Verified graph duals and disjoint reuse]
\label{prop:dual}
Orient any selected graph of pair constraints and let $M$ be its incidence
matrix, with one $+1$ and one $-1$ in each column. For a finite signed
multiplier vector $\lambda$, put $s=M\lambda$. Then
\begin{equation}
 D_G(\lambda)
 =s^\top y-\kappa\sum_{e=(i,j)}d(u_i,u_j)|\lambda_e|
       -\frac14\|s\|_2^2
 \ \leq\ n a_n.
 \label{eq:dual}
\end{equation}
Let $S_1,\ldots,S_K$ be disjoint subsets of cohort indices and let
$\underline A_k$ be verified nonnegative lower bounds on their minimum
squared residual sums. Then
\begin{equation}
 d_n=\frac1n\sum_k\underline A_k\leq a_n.
 \label{eq:reuse}
\end{equation}
Uncovered observations remain in the denominator $n$.
\end{proposition}

\begin{proof}
For a feasible vector $v$, the term
$\lambda^\top M^\top v-\kappa\sum_e d(u_i,u_j)|\lambda_e|$ is
nonpositive. Add it to $\|y-v\|_2^2$ and minimize over all $v$.
The minimizer is $y-s/2$, with value \eqref{eq:dual}, proving weak
duality without a convergence or optimality assumption. For
\eqref{eq:reuse}, restrict any full feasible vector to each disjoint
subset. Its residual sum on that subset is at least the corresponding
minimum; residuals at uncovered observations are nonnegative. Summing
and taking the full infimum proves the result.
\end{proof}

For equal blocks of size $b$ and
$\underline T_k=\underline A_k/(b-1)$, the reused value is
\begin{equation}
 d_n=\frac{K(b-1)}n\left(\frac1K\sum_k\underline T_k\right).
 \label{eq:normalization}
\end{equation}
This factor separates the centered block statistic from a full empirical
MSE. On disjoint vertex sets the graph-dual objectives add exactly.
Objectives from overlapping blocks cannot simply be summed because their
squared-norm terms have cross terms. They may instead be embedded in one
graph whose full objective is verified, or combined by a maximum of lower
values for the same $a_n$. Reuse requires no additional optimization solve;
it still incurs arithmetic and calibration work.

\subsection{Two lower calibrations with different independent sample sizes}

For $q,p\in[0,1]$, let
\[
 \operatorname{kl}(q\Vert p)
 =q\log(q/p)+(1-q)\log\{(1-q)/(1-p)\},
\]
with the usual zero-numerator and infinite boundary conventions.
For $c>0$, let $k_-(q,c)$ and $k_+(q,c)$ be the smallest and largest
$p$ with $\operatorname{kl}(q\Vert p)\leq c$. Both are nondecreasing
in $q$; increasing $c$ lowers $k_-$ and raises $k_+$.
The established bounded-mean Chernoff argument gives, for iid
$V_i\in[0,C]$, each of the one-sided statements
\begin{equation}
 C k_-(\bar V/C,\log(1/\alpha)/N)\leq\mathbb E V_1,
 \qquad
 \mathbb E V_1\leq C k_+(\bar V/C,\log(1/\alpha)/N)
 \label{eq:kl}
\end{equation}
with failure probability at most $\alpha$ \citep{hoeffding1963}.
The observations need not be Bernoulli. Indeed, convexity bounds their
exponential moments by $1-p+pe^t$, with $p=\mathbb EV_1/C$.
Independence and exponential Markov yield
$\exp\{-N\operatorname{kl}(q\Vert p)\}$ for the appropriate one-sided
empirical tail; inversion gives \eqref{eq:kl}. The two directions require
separate allocation when both are claimed. At $C=0$ use zero endpoints.

\begin{proposition}[An empirical-infimum lower endpoint]
\label{prop:empirical}
Let a fixed nonempty class have losses in $[0,B]$ and measurable empirical
infimum $a_n$ on $n$ iid observations. Any measurable returned value
$0\leq d_n\leq a_n$ pointwise satisfies
\begin{equation}
 B k_-(d_n/B,\log(1/\alpha)/n)\leq
       \inf_g\mathbb E\ell_g
 \label{eq:empirical}
\end{equation}
with probability at least $1-\alpha$.
\end{proposition}

\begin{proof}
Choose fixed population competitors $g_j$ whose risks $r_j$ decrease to
$r^*=\inf_g\mathbb E\ell_g$. They are proof devices, not predictors
chosen from calibration. Since
$d_n\leq a_n\leq n^{-1}\sum_i\ell_{g_j}(X_i,Y_i)$, monotonicity of
$k_-$ and \eqref{eq:kl} give
$\Pr\{B k_-(d_n/B,c)>r_j\}\leq\alpha$ for each $j$.
These events increase to $\{B k_-(d_n/B,c)>r^*\}$, where
$c=\log(1/\alpha)/n$. Continuity of probability proves the claim even
when the infimum is not attained.
\end{proof}

This fixed-competitor argument needs no uniform-complexity penalty for a
class \emph{lower} endpoint. It supplies neither a fitted-on-calibration
upper bound nor an excess-risk guarantee. It is an inherited calibration
ingredient, not a new general principle of statistical class falsification.

Now freeze a finite family of classes $m$ and fitted members $\widehat g_m$
using independent development data. Each record uses its declared known
range $B$, ordinary squared loss, target, and evaluation law. From $n$ iid
calibration observations, form $K_b=\lfloor n/b\rfloor\geq1$
predetermined disjoint blocks and pointwise lower scores
$0\leq\underline T_{m,b,k}\leq T_{m,b,k}$.
Also allow any pointwise value $0\leq d_{m,n}\leq a_{m,n}$ and verified
fitted-loss upper values $\ell_{m,i}\leq\overline\ell_{m,i}\leq B$.
For separately allocated probabilities, set
\begin{align}
 L^{\rm blk}_{m,b}
 &=C_b k_-\!\left(
     \frac{K_b^{-1}\sum_k\underline T_{m,b,k}}{C_b},
     \frac{\log(1/\alpha^{\rm blk}_{m,b})}{K_b}\right),
        \label{eq:Lblock}\\
 L^{\rm emp}_m
 &=B k_-(d_{m,n}/B,\log(1/\alpha^{\rm emp}_m)/n),
        \label{eq:Lemp}\\
 U_m&=B k_+\!\left(
       \frac{n^{-1}\sum_i\overline\ell_{m,i}}B,
       \frac{\log(1/\alpha^{\rm up}_m)}n\right).
        \label{eq:upper}
\end{align}
Take $L_m$ to be the maximum of the covered lower endpoints, or zero
when none are used. Numerical implementations report lower endpoints no
larger and upper endpoints no smaller than these exact-real quantities.

\begin{theorem}[Simultaneous class and fitted-risk certificate]
\label{thm:certificate}
If the finite event ledger is fixed before calibration and its allocated
failure probabilities sum to at most $\delta$, then, conditional on
development, with probability at least $1-\delta$,
\begin{equation}
 L_m\leq R_m^*\leq R_m(\widehat g_m)\leq U_m
 \label{eq:sandwich}
\end{equation}
simultaneously for every covered record. Each block endpoint separately
also lower-bounds its $\mathcal W_{m,b}$. Different events may reuse data.
\end{theorem}

\begin{proof}
For each block event, apply the lower inequality in \eqref{eq:kl} to
the exact iid scores $T_{m,b,k}\in[0,C_b]$. Their verified lower mean
cannot increase that endpoint. Theorem~\ref{thm:hierarchy} then gives
the class lower bound. Proposition~\ref{prop:empirical} gives the
empirical-infimum event. For a fitted upper, apply the upper inequality
in \eqref{eq:kl} to the exact iid losses; a larger verified mean is
conservative. Membership supplies $R_m^*\leq R_m(\widehat g_m)$.
A conditional union bound gives the common event; maxima of its lower
endpoints and outward numerical arithmetic preserve it. Integrating over
development yields the same unconditional probability statement.
\end{proof}

The block endpoint calibrates $K_b$ independent scores, not $n$ individual
observations. Larger blocks strengthen $\mathcal W_b$ but reduce $K_b$;
their finite-sample endpoints need not improve. The empirical-infimum
endpoint instead uses $n$ for the distinct pointwise objective $a_n$.
Neither calibration uniformly dominates the other. Numerical paths,
graphs, or partitions may be chosen adaptively when they still return a
pointwise lower value for the same fixed $a_n$. This does not authorize
changing a prediction class after inspecting calibration. Additional block
sizes, class choices, or sample-size looks require their own simultaneous
or sequential validity argument.

Independent episodes may contain arbitrarily dependent windows. An episode
average is one bounded fitted-loss observation. One independently sampled
origin per episode supplies the unconditional iid scalar observations for
\eqref{eq:block}. Counting overlapping windows as independent calibration
units is not justified by this theorem. Training extrema do not establish
an almost-sure target range, and capped-loss upper bounds cannot be combined
with the ordinary-MSE lower construction.

\subsection{Arithmetic semantics and the scope of a numerical certificate}

For \eqref{eq:dual}, treat retained multipliers and parameters according
to their declared exact semantics. An upper enclosure of each distance,
a lower enclosure of the linear term, and upper enclosures of subtracted
terms give a valid lower objective. The lower KL root and its final
rescaling must be rounded down; the upper root and rescaling must be
rounded up. Both directions use an upper enclosure of
$\log(1/\alpha)/N$. Rigorous deterministic arithmetic consumes no extra
statistical failure probability.

The declared feature map and the evaluated head must match the certificate.
A guarantee conditional on stored feature vectors does not certify their
upstream computation. Likewise, a mathematically Lipschitz affine head
and a rounded prediction program are distinct objects: rounding can
introduce discontinuities. Verifying output fidelity to a fixed head does
not establish that head's future prediction risk. These are separate
statistical and numerical claims. Coupled numerical enclosures transfer
safely through the monotone KL maps above; a variance-dependent empirical
Bernstein endpoint additionally needs a valid variance enclosure or iid
computed scores.

\section{Refinement and decisions within a declared family}
\label{sec:decisions}

Let $\Phi_j$ be fixed feature blocks and retain
$\Gamma_s=(\Phi_1,\ldots,\Phi_s)$ with metric
$\rho_s(u,v)=\max_{j\leq s}\|u_j-v_j\|_2$.
The scalar clipped affine family
\begin{equation}
 \mathcal F_{s,\kappa}
 =\left\{x\mapsto\operatorname{clip}_I\!\left(
 a+\sum_{j\leq s}\theta_j^\top\Phi_j(x)\right):
             \sum_{j\leq s}\|\theta_j\|_2\leq\kappa\right\}
 \label{eq:refinement}
\end{equation}
is a fitted subclass of the Lipschitz comparator. Coefficients and budgets
must use the physical target units after any inverse target transform.

\begin{proposition}[Retained blocks preserve the comparator]
\label{prop:refinement}
Every member of \eqref{eq:refinement} is $\kappa$-Lipschitz in
$\rho_s$. Appending an unchanged old-coordinate stack and a new block
embeds the earlier feasible family and cannot increase its optimal risk.
Appending an exact duplicate block leaves that predictor class and its
induced distance unchanged. The corresponding inclusions and duplication
invariance also hold for all Lipschitz predictors on the retained stack.
\end{proposition}

\begin{proof}
Cauchy--Schwarz and the triangle inequality bound the unprojected
prediction difference by
$\sum_j\|\theta_j\|_2\|u_j-v_j\|_2\leq\kappa\rho_s(u,v)$.
Clipping is nonexpansive. Assigning zero coefficients to a new block embeds
every old predictor. For a duplicate, merge its coefficients with those
of the copied block; the triangle inequality shows that the total norm
does not increase. This gives reverse inclusion. The maximum distance is
unchanged by an exact duplicate. For a general Lipschitz decoder,
composition with the coordinate projection embeds the old class because
that projection is nonexpansive. For duplication, composition with
$u\mapsto(u,u_j)$ gives the reverse inclusion because this map preserves
the maximum metric.
\end{proof}

Group norm constraints are established tools \citep{bach2008group}; here
they keep the sensitivity meaning compatible with retained blocks. These
are class-inclusion statements, not monotonicity claims for separately
trained estimators. A squared ridge penalty can change under duplication
even when the feasible class does not. Refitting old blocks or changing the
budget is a new comparison, not this particular nesting construction.

\begin{corollary}[Adequacy, exclusion, and minimum feasible cost]
\label{cor:decisions}
Fix an MSE tolerance $\tau\geq0$ and class costs $c_m\geq0$ before
calibration. On the joint event \eqref{eq:sandwich}, $U_m\leq\tau$
certifies an exhibited adequate predictor, while $L_m>\tau$ excludes
every decoder in that class. Otherwise the decision is unresolved.
Define a feasible class by the existence of a member with risk at most
$\tau$, and let $c^*$ be the minimum cost of a feasible class. With the
minimum of an empty set taken as $+\infty$,
\begin{equation}
 c_-:=\min_{m:L_m\leq\tau}c_m
 \ \leq\ c^*\ \leq\
 c_+:=\min_{m:U_m\leq\tau}c_m.
 \label{eq:cost}
\end{equation}
If $c_-=c_+<\infty$, an accepted fit attains the minimum feasible cost
within the declared family. If every class is excluded, that family is
infeasible at the requested tolerance.
\end{corollary}

\begin{proof}
Every accepted fitted member exhibits feasibility. Every genuinely feasible
class is nonexcluded, since its lower endpoint cannot exceed the risk of
an adequate member. These inclusions give the cost interval and both
conclusions without assuming attainment of a population infimum.
\end{proof}

An RMSE tolerance $\epsilon$ uses $\tau=\epsilon^2$. With several
required horizons, acceptance requires all covered upper tests, while one
covered lower test suffices to exclude a compatible joint class. When
$c_-<c_+$, the selected fit has minimum cost only among accepted fits.
Feature count, stored bytes, and runtime are different cost definitions;
the cost bracket does not establish savings in the total resources used
to search, fit, and certify candidates.

The comparison $L_c-U_r>0$ certifies that the exhibited replacement
outperforms every member of the coarse class by at least $L_c-U_r$.
Calling this a retained-coordinate refinement additionally requires the
class inclusion above. Its successful discovery is an empirical question.
Before attempting any adequacy test, the KL upper endpoint has the
deterministic zero-loss floor
\begin{equation}
 U\geq B(1-\alpha^{1/n}).
 \label{eq:floor}
\end{equation}
This follows from monotonicity and $k_+(0,c)=1-e^{-c}$. It is a necessary
feasibility check, not a power guarantee for a nonzero-risk predictor.
A certificate can be valid yet too weak to decide the requested tolerance.

\paragraph{Resolution beyond validity.}
Appendix~\ref{app:resolution} bounds finite-block population bias using an
internal uniform cover and gives a sufficient decision-resolution inequality.
It separates that bias from statistical calibration and verified numerical
slack. These sufficient bounds explain why a valid certificate can remain
inconclusive; they do not claim an optimal sample-complexity rate.
